% ============================================================================
% 02_architecture/02_positional_encoding.tex
% Section body for inclusion immediately after the corresponding \section
% command in the master book file. Follows section_body_template.tex.
% No document-level commands or package imports are included.
% ============================================================================

% ----------------------------------------------------------------------------
% Opening transition and conceptual motivation
% ----------------------------------------------------------------------------

Once tokenization has converted an input text into a sequence of discrete
tokens, the Transformer must construct numerical representations that retain
the information needed for subsequent computation. Token embeddings represent
the identities and learned content of the tokens, but a sequence is not
defined by token identity alone. The order in which tokens occur, and the
distances separating them, are also essential to linguistic structure. The
next architectural question is therefore how a Transformer can represent
sequence position and make that information available to its attention
mechanism.

This requirement follows from the computational organization of the
Transformer. Unlike recurrent architectures, which propagate information
through an explicitly ordered sequence of state updates, self-attention
computes interactions among token representations through query, key, and
value projections. This design enables parallel processing and direct
interactions between distant tokens, but the attention operation does not,
by itself, provide a complete representation of the positions occupied by
those tokens. Positional encoding supplies this missing structural
information. The distinction between token content and token position, and
the way in which they are combined, consequently form part of the mathematical
definition of a Transformer layer.

% ----------------------------------------------------------------------------
% Main exposition
% ----------------------------------------------------------------------------

\subsection{Why Self-Attention Requires Positional Information}
\label{subsec:pe-positional-necessity}

Let an input sequence contain \(L\) tokens, and let the embedding of the
token at position \(i\) be \(\mathbf{x}_i\in\mathbb{R}^{d_{\mathrm{model}}}\).
Collecting the embeddings as rows gives

\begin{equation}
    X =
    \begin{bmatrix}
        \mathbf{x}_1^{\mathsf T}\\
        \mathbf{x}_2^{\mathsf T}\\
        \vdots\\
        \mathbf{x}_L^{\mathsf T}
    \end{bmatrix}
    \in\mathbb{R}^{L\times d_{\mathrm{model}}}.
    \label{eq:pe-input-matrix}
\end{equation}

For one attention head, the query, key, and value matrices are obtained
through learned linear projections,

\begin{equation}
    Q=XW_Q,\qquad
    K=XW_K,\qquad
    V=XW_V,
    \label{eq:pe-qkv-projections}
\end{equation}

where \(W_Q,W_K\in\mathbb{R}^{d_{\mathrm{model}}\times d_k}\) and
\(W_V\in\mathbb{R}^{d_{\mathrm{model}}\times d_v}\). The query and key
dimension is \(d_k\), and the value dimension is \(d_v\). The scaled
dot-product attention scores and attention weights are

\begin{equation}
    S=\frac{QK^{\mathsf T}}{\sqrt{d_k}},
    \qquad
    A=\operatorname{softmax}_{\mathrm{row}}(S),
    \qquad
    Z=AV.
    \label{eq:pe-scaled-dot-product}
\end{equation}

The row-wise softmax normalizes each query's scores over the keys that it is
permitted to attend to. In unmasked self-attention, each row of \(A\) sums
to one. The factor \(\sqrt{d_k}\) controls the scale of the dot products as
the query and key dimension changes. Without this scaling, the variance of
the dot products can grow with \(d_k\), pushing softmax toward saturated
distributions and making optimization less well-conditioned. The formulation
in \cref{eq:pe-scaled-dot-product} is the self-attention basis from which the
positional mechanisms considered in this section are developed
\parencite{zhou2024rethinking}.

The limitation of content-only self-attention can be made precise through
permutation equivariance. Let \(P\in\mathbb{R}^{L\times L}\) be a permutation
matrix. Left multiplication by \(P\) rearranges the rows of a sequence
matrix. If the input is permuted, then

\begin{equation}
    X'=PX,\qquad Q'=PQ,\qquad K'=PK,\qquad V'=PV.
    \label{eq:pe-permuted-qkv}
\end{equation}

For unmasked self-attention, the score matrix becomes

\begin{align}
    S'
    &=\frac{Q'K'^{\mathsf T}}{\sqrt{d_k}}\\
    &=\frac{(PQ)(PK)^{\mathsf T}}{\sqrt{d_k}}\\
    &=P\frac{QK^{\mathsf T}}{\sqrt{d_k}}P^{\mathsf T}\\
    &=PSP^{\mathsf T}.
    \label{eq:pe-permuted-scores}
\end{align}

Because softmax is applied independently to each row, permuting both the
query rows and key columns permutes the attention matrix in the same way:
\(A'=PAP^{\mathsf T}\). The output is therefore

\begin{align}
    Z'
    &=A'V'\\
    &=PAP^{\mathsf T}PV\\
    &=PAV\\
    &=PZ.
    \label{eq:pe-permutation-equivariance}
\end{align}

Thus, content-only, unmasked self-attention is permutation-equivariant:
permuting the input representations permutes the outputs, but the operation
does not independently encode the original sequential positions. In an
autoregressive language model, a causal mask additionally restricts each
query to preceding positions. The mask establishes a direction of information
flow, but positional representations are still needed to provide the
attention computation with a more expressive account of absolute positions
and relative displacements. Positional encoding addresses this requirement
by making position-dependent information available to the network.

\subsection{Absolute Positional Encoding}
\label{subsec:pe-absolute}

An absolute positional encoding assigns a representation to each position
independently of the positions occupied by other tokens. If the token at
position \(p\) has content embedding \(\mathbf{x}_p\), its positional
representation is denoted by \(\mathbf{e}_p\). In the conventional additive
formulation, the input to the subsequent Transformer computation is

\begin{equation}
    \mathbf{h}_p=\mathbf{x}_p+\mathbf{e}_p,
    \qquad
    \mathbf{x}_p,\mathbf{e}_p
    \in\mathbb{R}^{d_{\mathrm{model}}}.
    \label{eq:pe-absolute-addition}
\end{equation}

The positional vector must have the same dimensionality as the token
embedding for this addition to be defined. Absolute positional representations
can be generated by predetermined functions or learned as model parameters.
The two alternatives considered by \textcite{zhou2024rethinking} are
sinusoidal positional encoding and learnable absolute positional encoding.

\paragraph{Sinusoidal positional encoding.}
Sinusoidal positional encoding constructs each positional vector from sine
and cosine functions with different frequencies. For an embedding dimension
\(d_{\mathrm{model}}\), the components are

\begin{align}
    \operatorname{PE}(p,2i)
    &=
    \sin\left(
        \frac{p}{10000^{2i/d_{\mathrm{model}}}}
    \right),
    \label{eq:pe-sinusoidal-even}\\
    \operatorname{PE}(p,2i+1)
    &=
    \cos\left(
        \frac{p}{10000^{2i/d_{\mathrm{model}}}}
    \right),
    \label{eq:pe-sinusoidal-odd}
\end{align}

where \(p\) is the position index and \(i\) indexes a pair of embedding
dimensions. Defining the angular frequency

\begin{equation}
    \omega_i=10000^{-2i/d_{\mathrm{model}}},
    \label{eq:pe-sinusoidal-frequency}
\end{equation}

the two components associated with frequency \(\omega_i\) can be written as

\begin{equation}
    \mathbf{p}_i(p)=
    \begin{bmatrix}
        \sin(p\omega_i)\\
        \cos(p\omega_i)
    \end{bmatrix}.
    \label{eq:pe-sinusoidal-pair}
\end{equation}

The frequencies decrease geometrically as the pair index increases. The
lower-index pairs therefore change more rapidly with position, while the
higher-index pairs change more slowly. The complete positional vector
combines these frequency scales, giving the model a deterministic
multiscale representation of position.

For example, when \(d_{\mathrm{model}}=4\), there are two frequency pairs.
Their frequencies are

\begin{equation}
    \omega_0=10000^0=1,
    \qquad
    \omega_1=10000^{-1/2}=0.01.
    \label{eq:pe-sinusoidal-numeric-frequencies}
\end{equation}

At position \(p=2\), the positional vector is

\begin{align}
    \operatorname{PE}(2)
    &=
    \begin{bmatrix}
        \sin(2)\\
        \cos(2)\\
        \sin(0.02)\\
        \cos(0.02)
    \end{bmatrix}\\
    &\approx
    \begin{bmatrix}
        0.9093\\
        -0.4161\\
        0.0200\\
        0.9998
    \end{bmatrix}.
    \label{eq:pe-sinusoidal-numeric-example}
\end{align}

The first pair varies substantially over a small positional change, whereas
the second pair varies slowly. Their combination allows positions to be
represented through multiple periodic scales rather than through a single
monotonically increasing coordinate.

A useful property of this construction is that a fixed positional
displacement can be represented by a linear transformation. For a fixed
offset \(k\), the trigonometric addition identities give

\begin{equation}
    \mathbf{p}_i(p+k)=
    \begin{bmatrix}
        \cos(k\omega_i)&\sin(k\omega_i)\\
        -\sin(k\omega_i)&\cos(k\omega_i)
    \end{bmatrix}
    \mathbf{p}_i(p).
    \label{eq:pe-sinusoidal-relative-linear}
\end{equation}

The transformation in \cref{eq:pe-sinusoidal-relative-linear} depends on
the offset \(k\), not on the absolute position \(p\). Although the encoding
is defined using absolute positions, this property gives a model a
mathematical route to learning relationships associated with relative
displacements. The sinusoidal construction is deterministic and does not
require a separately learned positional parameter vector for every position
\parencite{zhou2024rethinking}.

\paragraph{Learnable absolute positional encoding.}
An alternative is to learn a positional vector for each position. For a
maximum represented sequence length \(L_{\max}\), the positional parameters
can be arranged as

\begin{equation}
    E_{\mathrm{pos}}\in
    \mathbb{R}^{L_{\max}\times d_{\mathrm{model}}},
    \qquad
    \mathbf{e}_p=E_{\mathrm{pos}}[p,:].
    \label{eq:pe-learned-table}
\end{equation}

The token and positional vectors are then combined using
\cref{eq:pe-absolute-addition}. The number of positional parameters is

\begin{equation}
    N_{\mathrm{pos}}=
    L_{\max}d_{\mathrm{model}}.
    \label{eq:pe-learned-parameter-count}
\end{equation}

Unlike the sinusoidal method, the learned table does not impose a
predetermined frequency structure across positions. Its positional vectors
are optimized with the other model parameters. This flexibility makes the
representation data-dependent, while the explicit table also ties the
learned parameters to the positions represented during training. The two
methods therefore differ in how positional structure is supplied: sinusoidal
encoding imposes a fixed multiscale construction, whereas learnable encoding
estimates the positional vectors through optimization.

\subsection{Combining Content and Position: Addition and Concatenation}
\label{subsec:pe-addition-concatenation}

The construction of a positional vector is only one part of the design.
The model must also determine how positional information is combined with
token content before the attention projections are applied. Addition is the
conventional method, but concatenation provides an alternative that keeps
the two components in distinct parts of the representation.

Under addition, the query and key vectors at positions \(i\) and \(j\) are

\begin{equation}
    \mathbf{q}_i=(\mathbf{x}_i+\mathbf{p}_i)W_Q,
    \qquad
    \mathbf{k}_j=(\mathbf{x}_j+\mathbf{p}_j)W_K.
    \label{eq:pe-additive-qk}
\end{equation}

Let \(W_{QK}=W_QW_K^{\mathsf T}\). The unnormalized dot product can be
expanded into four interactions:

\begin{align}
    \mathbf{q}_i\mathbf{k}_j^{\mathsf T}
    &=
    (\mathbf{x}_i+\mathbf{p}_i)
    W_{QK}
    (\mathbf{x}_j+\mathbf{p}_j)^{\mathsf T}
    \nonumber\\
    &=
    \mathbf{x}_iW_{QK}\mathbf{x}_j^{\mathsf T}
    +\mathbf{p}_iW_{QK}\mathbf{x}_j^{\mathsf T}
    \nonumber\\
    &\quad+
    \mathbf{x}_iW_{QK}\mathbf{p}_j^{\mathsf T}
    +\mathbf{p}_iW_{QK}\mathbf{p}_j^{\mathsf T}.
    \label{eq:pe-additive-four-terms}
\end{align}

The first term is a content-to-content interaction. The second is a
position-to-content interaction, the third is a content-to-position
interaction, and the fourth is a position-to-position interaction. In the
additive formulation, these interactions share the same learned
query-key interaction matrix \(W_{QK}\). They are consequently not
parameterized as four independent transformations.

\textcite{zhou2024rethinking} use this decomposition to motivate a
concatenation-based alternative. To maintain the same total input dimension,
the content and positional components are each assigned half of the
representation dimension:

\begin{equation}
    \mathbf{x}'_i,\mathbf{p}'_i
    \in\mathbb{R}^{d_{\mathrm{model}}/2},
    \qquad
    \mathbf{h}'_i=
    \mathbf{x}'_i\oplus\mathbf{p}'_i
    \in\mathbb{R}^{d_{\mathrm{model}}},
    \label{eq:pe-concatenated-input}
\end{equation}

where \(\oplus\) denotes vector concatenation. The concatenated
representation is projected into queries and keys in the usual way. If the
resulting query-key interaction matrix is partitioned into four blocks,

\begin{equation}
    W_{QK}=
    \begin{bmatrix}
        W_{cc}&W_{cp}\\
        W_{pc}&W_{pp}
    \end{bmatrix},
    \label{eq:pe-concat-block-matrix}
\end{equation}

then the score contains separately parameterized content-to-content,
content-to-position, position-to-content, and position-to-position
interactions:

\begin{align}
    \mathbf{q}'_i\mathbf{k}'_j{}^{\mathsf T}
    &=
    \mathbf{x}'_iW_{cc}\mathbf{x}'_j{}^{\mathsf T}
    +\mathbf{x}'_iW_{cp}\mathbf{p}'_j{}^{\mathsf T}
    \nonumber\\
    &\quad+
    \mathbf{p}'_iW_{pc}\mathbf{x}'_j{}^{\mathsf T}
    +\mathbf{p}'_iW_{pp}\mathbf{p}'_j{}^{\mathsf T}.
    \label{eq:pe-concat-four-terms}
\end{align}

The block decomposition makes the distinction between the two methods
explicit. Addition uses a shared interaction matrix for the four terms,
whereas concatenation allows the corresponding blocks to be learned
independently. The dimension reduction in \cref{eq:pe-concatenated-input}
keeps the concatenated input dimension equal to that of the additive
representation. The conceptual difference is illustrated in
\cref{fig:pe-addition-concatenation}.

Zhou et al. argue that content and positional vectors represent distinct
types of information and that independently parameterized interaction
spaces can therefore be beneficial. Their experiments compare addition and
concatenation in several single-scale and multiscale Vision Transformer
models. They report performance differences ranging from approximately
\(0.1\) to \(0.5\) percentage points in top-1 ImageNet-1K accuracy in
favor of concatenation for several evaluated models, without a systematic
increase in computation overhead. These are results from vision
architectures and image-classification experiments; they should not be
interpreted as direct evidence that concatenation is superior for language
models. The contribution relevant to language-model architecture is the
analysis of how the choice of combination operation changes the
parameterization of content-position interactions \parencite{zhou2024rethinking}.

\subsection{Relative Positional Encoding and Rotary Position Embedding}
\label{subsec:pe-relative-rope}

Absolute positional encoding assigns a representation to each individual
position. Many relationships between tokens, however, are naturally
described by the displacement between their positions. If two tokens occur
at positions \(m\) and \(n\), their signed relative displacement is

\begin{equation}
    \Delta=m-n.
    \label{eq:pe-relative-displacement}
\end{equation}

For example, the pairs \((m,n)=(10,7)\) and \((m,n)=(100,97)\) have
different absolute positions but the same displacement, \(\Delta=3\).
A relative positional method aims to make the interaction between the
tokens depend on this relationship. One family of relative methods uses
learnable representations indexed by relative displacement. Rotary
Position Embedding (RoPE) achieves a related objective by rotating query
and key representations according to their absolute positions so that
their inner product depends on the relative displacement
\parencite{zhou2024rethinking,zhang2024foundmiddle}.

\paragraph{Two-dimensional rotations.}
The mathematical basis of RoPE is a rotation in a two-dimensional real
vector space. For an angle \(\phi\), define

\begin{equation}
    R(\phi)=
    \begin{bmatrix}
        \cos\phi&-\sin\phi\\
        \sin\phi&\cos\phi
    \end{bmatrix}.
    \label{eq:pe-rotation-matrix}
\end{equation}

For a vector \(\mathbf{x}\in\mathbb{R}^2\), the rotated vector is
\(R(\phi)\mathbf{x}\). The matrix is orthogonal:

\begin{equation}
    R(\phi)^{\mathsf T}R(\phi)=I,
    \qquad
    R(\phi)^{\mathsf T}=R(-\phi).
    \label{eq:pe-rotation-orthogonality}
\end{equation}

It follows that rotations preserve the Euclidean norm,

\begin{equation}
    \|R(\phi)\mathbf{x}\|_2^2
    =
    \mathbf{x}^{\mathsf T}R(\phi)^{\mathsf T}
    R(\phi)\mathbf{x}
    =
    \|\mathbf{x}\|_2^2.
    \label{eq:pe-rotation-norm}
\end{equation}

The composition of two rotations adds their angles:

\begin{equation}
    R(\alpha)R(\beta)=R(\alpha+\beta).
    \label{eq:pe-rotation-composition}
\end{equation}

These identities are the essential algebraic properties used by RoPE.
Equivalently, a two-dimensional vector can be represented as a complex
number \(z=x_1+\mathrm{i}x_2\), and multiplication by
\(e^{\mathrm{i}\phi}\) performs the same rotation through Euler's identity.

\paragraph{Frequency-dependent rotations.}
Let the query and key dimension of one attention head be \(d_h\), assumed
even. RoPE groups the vector coordinates into \(d_h/2\) pairs and assigns
a different angular frequency to each pair. With the frequency convention
used in the papers considered here,

\begin{equation}
    \theta_j=10000^{-2j/d_h},
    \qquad
    j=0,1,\ldots,\frac{d_h}{2}-1.
    \label{eq:pe-rope-frequencies}
\end{equation}

For position \(m\), the complete rotation is a block-diagonal matrix,

\begin{equation}
    R_m=
    \operatorname{diag}
    \left(
        R(m\theta_0),
        R(m\theta_1),
        \ldots,
        R(m\theta_{d_h/2-1})
    \right).
    \label{eq:pe-rope-block-rotation}
\end{equation}

The transformed query and key vectors are

\begin{equation}
    \widetilde{\mathbf{q}}_m=R_m\mathbf{q}_m,
    \qquad
    \widetilde{\mathbf{k}}_n=R_n\mathbf{k}_n.
    \label{eq:pe-rope-qk}
\end{equation}

The query-key inner product is then

\begin{align}
    \widetilde{\mathbf{q}}_m^{\mathsf T}
    \widetilde{\mathbf{k}}_n
    &=
    (R_m\mathbf{q}_m)^{\mathsf T}
    (R_n\mathbf{k}_n)\\
    &=
    \mathbf{q}_m^{\mathsf T}
    R_m^{\mathsf T}R_n
    \mathbf{k}_n\\
    &=
    \mathbf{q}_m^{\mathsf T}
    R_{n-m}
    \mathbf{k}_n.
    \label{eq:pe-rope-relative-inner-product}
\end{align}

The last equality follows from the orthogonality and composition
properties in \cref{eq:pe-rotation-orthogonality,eq:pe-rotation-composition}.
It establishes the central property of RoPE: although each query and key
is transformed using its own absolute position, their attention
interaction contains a rotation determined by the relative displacement
\(n-m\). The sign of the displacement depends on the convention used for
the rotation matrix; the essential property is dependence on the difference
between positions. The structure is summarized in
\cref{fig:pe-rope-relative-position}.

The multiple frequencies in \cref{eq:pe-rope-frequencies} cause different
coordinate pairs to rotate at different rates. Higher-frequency components
change phase more rapidly as position changes, while lower-frequency
components change more gradually. RoPE therefore introduces positional
information at multiple frequency scales directly into query-key
interactions. It does not require a separate learned positional embedding
table, because the transformations are determined by the position indices
and the frequency schedule \parencite{zhang2024foundmiddle}.

\subsection{Long-Context Position Rescaling and Multi-Scale Positional Encoding}
\label{subsec:pe-mspoe}

The ability to encode relative displacement does not guarantee that a
language model will use all positions in a long input equally well.
\textcite{zhang2024foundmiddle} study the ``lost-in-the-middle''
phenomenon: relevant information can be overlooked when it occurs in the
middle of a long context, even when the model can process the complete
sequence. They identify two contributing factors. First, causal attention
can favor initial tokens because these tokens are available to a larger
number of subsequent queries. Second, the long-distance behavior associated
with RoPE can reduce attention to distant tokens, including tokens that
contain relevant information.

The latter observation should be interpreted as an empirical property of
the resulting attention behavior, not as a theorem that every individual
RoPE query-key score decreases monotonically with distance. The rotation
matrices themselves preserve vector norms, and their inner-product
contributions vary with relative phase. The paper's motivation is that the
combined behavior of RoPE-based attention in the evaluated models can
contribute to reduced use of distant context.

\paragraph{Uniform positional rescaling.}
Since RoPE uses position indices to determine rotation angles, the
effective positional distances can be changed by rescaling the indices.
For a positive scaling factor \(r\), replace position \(m\) by

\begin{equation}
    m'=\frac{m}{r}.
    \label{eq:pe-position-rescaling}
\end{equation}

For two positions \(m\) and \(n\), the rescaled displacement becomes

\begin{equation}
    \Delta'
    =
    \frac{m}{r}-\frac{n}{r}
    =
    \frac{m-n}{r}
    =
    \frac{\Delta}{r}.
    \label{eq:pe-rescaled-displacement}
\end{equation}

When \(r>1\), the effective positional distance is compressed. When
\(0<r<1\), it is expanded, and \(r=1\) recovers the original RoPE
positions. For example, two tokens separated by \(3{,}000\) positions
have an effective separation of

\begin{equation}
    \Delta'=\frac{3000}{1.5}=2000
    \label{eq:pe-rescaling-numeric-example}
\end{equation}

when \(r=1.5\). The sequence itself is not shortened and no tokens are
removed; only the positional indices used by the rotary transformation
are remapped. This can reduce the phase difference between distant
positions and potentially alleviate the long-distance decay observed by
the authors.

Uniform rescaling has a trade-off. A factor that is too large can alter
the positional patterns learned during pre-training and can change which
regions of the context the model favors. Fractional rescaled positions
may also differ from the integer position indices encountered during
pre-training, creating a positional out-of-distribution concern. In the
experiments reported by Zhang et al., small scaling factors such as
\(0.5\) favor more recent context, whereas larger factors such as \(2.5\)
can favor the beginning of the context. A factor near \(1.5\) improves
average performance in their initial rescaling experiments, but a single
factor does not account for the different positional behavior of
individual attention heads.

\paragraph{Position-awareness across attention heads.}
The motivation for Multi-Scale Positional Encoding (Ms-PoE) is the
observation that attention heads differ in their ability to follow relevant
information when its position changes. Some heads shift their attention
toward the relevant content as that content moves through the input.
Other heads remain more strongly concentrated on the beginning and end
regions. Zhang et al. use the term \emph{position-aware} for heads that
can capture relevant information across these positional changes.

For a layer \(k\), let the attention distribution of head \(h\) for the
most recent query be

\begin{equation}
    \mathbf{A}_{k,h}
    =
    [A_{k,h,1},A_{k,h,2},\ldots,A_{k,h,L}],
    \qquad
    \sum_{t=1}^{L}A_{k,h,t}=1.
    \label{eq:pe-mspoe-attention-vector}
\end{equation}

Its mean component is \(1/L\). The Position-Awareness Score is defined as

\begin{equation}
    S^P_{k,h}
    =
    \frac{1}{L}
    \sum_{t=1}^{L}
    \mathbb{I}
    \left(
        A_{k,h,t}
        \geq
        \alpha
        \frac{1}{L}
        \sum_{j=1}^{L}A_{k,h,j}
    \right),
    \label{eq:pe-position-awareness-score}
\end{equation}

where \(\mathbb{I}(\cdot)\) is the indicator function and \(\alpha\) is
a threshold hyperparameter. The experiments use \(\alpha=3\). Because
the attention weights sum to one, the threshold in this case is
\(3/L\). The score is the fraction of positions whose attention weights
are at least three times the mean attention weight. A larger value
indicates that a larger fraction of positions receive attention above
this threshold. This score is an operational measure introduced by the
authors; it is not itself a direct measure of semantic understanding or
retrieval accuracy.

For each layer, heads are sorted in descending order of their
Position-Awareness Scores. The most position-aware head receives a
rescaling factor near one, while heads with lower scores receive larger
factors. The authors use a linearly spaced set of factors from \(1.2\)
to \(1.8\). For the head at sorted rank \(i\), the assigned factor is

\begin{equation}
    r_i=
    1.2+
    \frac{(i-1)(1.8-1.2)}{n_h-1},
    \qquad i=1,\ldots,n_h,
    \label{eq:pe-mspoe-scaling-ratio}
\end{equation}

where \(n_h\) is the number of attention heads in the layer. Thus, the
first head in the descending positional-awareness ordering receives
\(r_1=1.2\), and the last receives \(r_{n_h}=1.8\). The factors are
assigned separately within each layer. The resulting head-specific
positional transformation is

\begin{equation}
    \widetilde{\mathbf{q}}_{k,h,m}
    =
    R\left(\frac{m}{r_{k,h}}\right)
    \mathbf{q}_{k,h,m},
    \qquad
    \widetilde{\mathbf{k}}_{k,h,n}
    =
    R\left(\frac{n}{r_{k,h}}\right)
    \mathbf{k}_{k,h,n},
    \label{eq:pe-mspoe-headwise-rope}
\end{equation}

where \(R(\cdot)\) denotes the frequency-dependent block rotation
defined earlier. The relative displacement used by head \(h\) in layer
\(k\) is consequently \((m-n)/r_{k,h}\). Different heads therefore
operate with different effective positional scales. The arrangement is
illustrated in \cref{fig:pe-mspoe-headwise}.

\paragraph{Inference procedure and computational considerations.}
Ms-PoE is applied to a pre-trained model that already uses RoPE. During
prefilling, the model obtains the attention distribution of the most
recent query for every head in every layer. It computes the
Position-Awareness Scores, sorts the heads within each layer, and assigns
the linearly spaced rescaling factors. These factors remain fixed during
the subsequent decoding stage. The procedure does not require additional
fine-tuning or additional trainable positional parameters
\parencite{zhang2024foundmiddle}.

\begin{algorithm}[H]
    \caption{Head-wise positional rescaling in Ms-PoE}
    \label{alg:pe-mspoe}
    \KwIn{Pre-trained RoPE-based language model; input sequence; number
    of layers \(n_l\); number of heads \(n_h\); threshold \(\alpha\);
    scaling interval \([1.2,1.8]\).}
    \KwOut{Generated sequence using head-specific positional rescaling.}
    \For{\(k=1\) \KwTo \(n_l\)}{
        Obtain the most recent query and the keys for layer \(k\)\;
        \For{\(h=1\) \KwTo \(n_h\)}{
            Obtain the attention distribution
            \(\mathbf{A}_{k,h}\)\;
            Compute \(S^P_{k,h}\) using
            \cref{eq:pe-position-awareness-score}\;
        }
        Sort the heads by \(S^P_{k,h}\) in descending order\;
        Assign the factors \(r_i\) using
        \cref{eq:pe-mspoe-scaling-ratio}\;
        Associate each assigned factor with its sorted head\;
    }
    Replace the original RoPE position indices by the corresponding
    head-specific rescaled indices during decoding\;
    Generate tokens autoregressively using the fixed factors\;
\end{algorithm}

The attention vectors used for the score are already part of the
prefilling attention computation. Given these vectors, evaluating the
threshold indicator for all heads and positions requires
\(O(n_l n_h L)\) comparisons. Sorting the \(n_h\) heads in each layer
requires \(O(n_l n_h\log n_h)\) comparisons. Storing the scores and
assigned factors requires \(O(n_l n_h)\) scalar values. These costs
describe the additional score-reduction and ordering operations implied
by the procedure; the method does not introduce a new trainable
parameter matrix or a new attention layer. Once the factors are fixed,
the decoding stage uses the modified RoPE transformations for the
existing heads.

\paragraph{Experimental evaluation.}
Zhang et al. evaluate Ms-PoE on long-context generation and retrieval
tasks using Llama-2, StableBeluga, and Vicuna models. On the
ZeroSCROLLS benchmark, they report average improvements over the
corresponding baselines for the evaluated models. The reported average
gain is \(1.6\) points for Llama-2-7B-Chat, \(0.9\) for
Llama-2-13B-Chat, \(2.6\) for StableBeluga-7B, \(0.6\) for
StableBeluga-13B, \(3.8\) for Vicuna-7B, and \(3.3\) for
Vicuna-7B-16K. The results indicate improvements across the reported
model sizes and context configurations, although the magnitude varies
by model and task \parencite{zhang2024foundmiddle}.

The multi-document question-answering and key-value retrieval
experiments vary the location of the relevant document or key-value
pair. For Vicuna-7B, the average multi-document question-answering
accuracy increases from \(61.12\%\) for the baseline to \(65.04\%\)
with Ms-PoE. For key-value retrieval, the corresponding average
increases from \(36.24\%\) to \(79.96\%\). The authors average results
over \(500\) samples and use the difference between the best and worst
accuracy across relevant-information positions as a measure of
positional sensitivity:

\begin{equation}
    \operatorname{GapAccuracy}
    =
    \max_{p}\operatorname{Acc}(p)
    -
    \min_{p}\operatorname{Acc}(p),
    \label{eq:pe-gap-accuracy}
\end{equation}

where \(\operatorname{Acc}(p)\) is the accuracy when the relevant
information is placed at position \(p\). A smaller gap indicates less
variation in accuracy as the location of relevant information changes.
The results show that Ms-PoE improves average retrieval performance
and reduces the positional gap in the evaluated settings
\parencite{zhang2024foundmiddle}.

The ablation experiments also examine the way scaling factors are
assigned. In the reported multi-document question-answering experiment,
the position-awareness ordering achieves an average accuracy of
\(65.3\%\), compared with \(61.7\%\) for random assignment,
\(57.8\%\) for sequential assignment, and \(62.3\%\) for
attention-entropy-based assignment. The baseline average is
\(62.1\%\). These results support the authors' claim that head-wise
positional-awareness information provides a more effective assignment
criterion than the tested alternatives. They do not establish that the
same ordering criterion is optimal for every architecture or task.

\subsection{Limitations and Connection to the Next Architectural Components}
\label{subsec:pe-limitations-transition}

Positional encoding determines how a Transformer represents token
locations and how positional information enters attention interactions.
The methods considered here make different design choices. Sinusoidal
encoding imposes a fixed multiscale trigonometric structure; learnable
absolute encoding estimates a positional table during training; addition
and concatenation determine how content and position are combined;
RoPE uses position-dependent rotations whose query-key interaction
depends on relative displacement; and Ms-PoE modifies RoPE at inference
by assigning different rescaling factors to different attention heads.

Each construction has limitations. A learned absolute positional table
is tied to the positions represented by its parameters. Sinusoidal
encoding supplies a fixed frequency structure, but the model must still
learn to use that structure through its projections. Concatenation
changes the parameterization of content-position interactions, but the
experimental comparison reported by Zhou et al. concerns vision
architectures and cannot by itself establish an advantage for language
models. RoPE provides a relative-position interaction through rotations,
but its behavior over long distances can still be associated with
reduced utilization of distant context. Ms-PoE addresses this issue in
the evaluated models, but its head-wise score is a threshold-based
proxy, its factors are selected from a specified interval, and
rescaling can create positional patterns that differ from those seen
during pre-training.

The central architectural conclusion is that positional encoding is
not merely an auxiliary vector added to token embeddings. It determines
how sequence order and relative displacement influence the attention
computation, and therefore affects how a language model can represent
and use information distributed across a sequence. With token content
and positional relationships now specified as inputs to the attention
mechanism, the next architectural question is how each token
representation is further transformed after attention. This leads to
the feed-forward component of the Transformer, which applies
position-wise nonlinear transformations to the representations
constructed by attention.

% ----------------------------------------------------------------------------
% Figures
% ----------------------------------------------------------------------------
% Figure source code is stored separately in the root-level figures/
% directory and is included here. Alt text is recorded in comments.

% Figure source: addition and concatenation of token and positional vectors.
% Uses the standard LaTeX picture environment; no additional package is needed.

\begin{figure}[htbp]
    \centering
    \setlength{\unitlength}{1mm}
    \resizebox{\textwidth}{!}{%
    \begin{picture}(174,64)
        \put(0,59){\makebox(0,0)[l]{\textbf{Addition}}}
        \put(0,42){\framebox(29,10){\shortstack{Token embedding\\$\mathbf{x}_i\in\mathbb{R}^{d}$}}}
        \put(0,15){\framebox(29,10){\shortstack{Position vector\\$\mathbf{p}_i\in\mathbb{R}^{d}$}}}
        \put(29,47){\line(1,-1){11}}
        \put(29,20){\line(1,1){11}}
        \put(43,33){\makebox(0,0){\Large \(+\)}}
        \put(47,33){\vector(1,0){11}}
        \put(58,27){\framebox(35,12){\shortstack{Combined vector\\$\mathbf{x}_i+\mathbf{p}_i$}}}

        \put(104,59){\makebox(0,0)[l]{\textbf{Concatenation}}}
        \put(104,42){\framebox(27,10){\shortstack{Content\\$\mathbf{x}'_i$}}}
        \put(104,15){\framebox(27,10){\shortstack{Position\\$\mathbf{p}'_i$}}}
        \put(131,47){\line(1,-1){10}}
        \put(131,20){\line(1,1){10}}
        \put(144,33){\makebox(0,0){\Large \(\oplus\)}}
        \put(148,33){\vector(1,0){8}}
        \put(156,27){\framebox(18,12){\shortstack{Joined\\vector}}}
    \end{picture}%
    }
    \caption{Two ways of incorporating positional information into token
    representations. Addition superimposes full-dimensional vectors,
    whereas concatenation joins half-dimensional content and positional
    components while retaining the same total dimension. Adapted
    conceptually from the comparison in \textcite{zhou2024rethinking}.}
    \label{fig:pe-addition-concatenation}
\end{figure}

% Alt text: Two parallel diagrams compare adding a full-dimensional
% positional vector to a token embedding with concatenating two
% half-dimensional content and positional vectors.

% Figure source: positional rotations in RoPE.
% Uses the standard LaTeX picture environment; no additional package is needed.

\begin{figure}[htbp]
    \centering
    \setlength{\unitlength}{1mm}
    \resizebox{\textwidth}{!}{%
    \begin{picture}(158,66)
        \put(0,53){\framebox(28,10){Query $\mathbf{q}_m$}}
        \put(33,58){\vector(1,0){9}}
        \put(42,53){\framebox(36,10){$R_m=R(m\boldsymbol{\theta})$}}
        \put(80,58){\vector(1,0){9}}
        \put(89,53){\framebox(34,10){$\widetilde{\mathbf{q}}_m$}}

        \put(0,31){\framebox(28,10){Key $\mathbf{k}_n$}}
        \put(33,36){\vector(1,0){9}}
        \put(42,31){\framebox(36,10){$R_n=R(n\boldsymbol{\theta})$}}
        \put(80,36){\vector(1,0){9}}
        \put(89,31){\framebox(34,10){$\widetilde{\mathbf{k}}_n$}}

        \put(0,10){\makebox(0,0)[l]{$
        \widetilde{\mathbf{q}}_m^{\mathsf T}
        \widetilde{\mathbf{k}}_n
        =
        \mathbf{q}_m^{\mathsf T}R_{n-m}\mathbf{k}_n
        $}}
    \end{picture}%
    }
    \caption{Rotary Position Embedding applies position-dependent rotations
    to queries and keys. Their inner product contains a rotation determined
    by the relative displacement between positions. Schematic based on
    \textcite{zhang2024foundmiddle}.}
    \label{fig:pe-rope-relative-position}
\end{figure}

% Alt text: Query and key vectors at positions m and n are separately
% rotated by their position-dependent RoPE transformations; their inner
% product is expressed using a rotation determined by the relative offset.

% Figure source: head-wise rescaling in Multi-Scale Positional Encoding.
% Uses the standard LaTeX picture environment; no additional package is needed.

\begin{figure}[htbp]
    \centering
    \setlength{\unitlength}{1mm}
    \resizebox{\textwidth}{!}{%
    \begin{picture}(166,72)
        \put(1,64){\makebox(0,0)[l]{\textbf{Attention heads and positional-awareness scores}}}

        \put(1,48){\framebox(37,10){Head A: high $S^P$}}
        \put(1,32){\framebox(37,10){Head B: medium $S^P$}}
        \put(1,16){\framebox(37,10){Head C: low $S^P$}}

        \put(38,53){\vector(1,0){12}}
        \put(38,37){\vector(1,0){12}}
        \put(38,21){\vector(1,0){12}}

        \put(50,48){\framebox(31,10){$r=1.2$}}
        \put(50,32){\framebox(31,10){$r=1.5$}}
        \put(50,16){\framebox(31,10){$r=1.8$}}

        \put(81,53){\vector(1,0){10}}
        \put(81,37){\vector(1,0){10}}
        \put(81,21){\vector(1,0){10}}

        \put(91,48){\framebox(35,10){$m/r_h$ in RoPE}}
        \put(91,32){\framebox(35,10){$m/r_h$ in RoPE}}
        \put(91,16){\framebox(35,10){$m/r_h$ in RoPE}}

        \put(126,53){\line(1,-1){8}}
        \put(126,37){\line(1,0){8}}
        \put(126,21){\line(1,1){8}}
        \put(136,37){\vector(1,0){8}}
        \put(144,31){\framebox(22,12){\shortstack{Multi-scale\\attention}}}

        \put(1,5){\makebox(0,0)[l]{\footnotesize Higher positional awareness receives a factor closer to one.}}
    \end{picture}%
    }
    \caption{Conceptual flow of Ms-PoE. Heads are ranked by their
    position-awareness scores and receive different rescaling factors;
    higher-scoring heads are changed less, while lower-scoring heads
    receive stronger positional compression. The values shown are
    illustrative points in the interval used by \textcite{zhang2024foundmiddle}.}
    \label{fig:pe-mspoe-headwise}
\end{figure}

% Alt text: Attention heads are ranked by positional-awareness scores
% and assigned different rescaling factors between 1.2 and 1.8 before
% head-specific RoPE transformations are applied.
